\documentclass[11pt]{article}
\usepackage[a4paper,margin=25mm]{geometry}
\usepackage[T1]{fontenc}
\usepackage{lmodern,amsmath,amssymb,amsthm,booktabs,array}
\usepackage[hidelinks]{hyperref}
\usepackage{xurl}
\newtheorem{theorem}{Theorem}
\newtheorem{lemma}[theorem]{Lemma}
\newtheorem{corollary}[theorem]{Corollary}
\newcommand{\Sep}{\mathcal S}
\newcommand{\ket}[1]{\lvert #1\rangle}
\newcommand{\bra}[1]{\langle #1\rvert}
\newcommand{\proj}[1]{\ket{#1}\!\bra{#1}}
\newcommand{\Tr}{\operatorname{Tr}}
\newcommand{\diag}{\operatorname{diag}}
\newcommand{\Rea}{\operatorname{Re}}
\newcommand{\Ima}{\operatorname{Im}}
\newcommand{\E}{E_{\mathrm{HS}}}
\newcommand{\D}{D_{\mathrm{HS}}}
\hypersetup{pdftitle={Refuting the X-MEMS closest-separable-state conjecture in arXiv:2209.01430v2, Appendix A, Eq. (A11)},pdfauthor={Haobo Ma and Wenlin Zhang}}
\title{Refuting the X-MEMS closest-separable-state conjecture in\break arXiv:2209.01430v2, Appendix A, Eq.~(A11)}
\author{Haobo Ma\thanks{ChronoAI Pte Ltd, Singapore.}\and
Wenlin Zhang\thanks{National University of Singapore; ChronoAI Pte Ltd, Singapore.}}
\date{2 November 2026}
\begin{document}
\maketitle
\begin{abstract}
We refute the closest-fully-separable-state shape conjecture of Consiglio,
Apollaro and Wie\'sniak, arXiv:2209.01430v2, Appendix A, Eq.~(A11), already
for three qubits on the low branch at $\gamma=11/112$.
Every fully separable state of the proposed shape has squared
Hilbert--Schmidt distance at least $109/23520$, whereas an explicit mixture
of ten pure product states has distance
$417965512477/90316800000000$.
The counterexample and exact restricted minimum have self-contained written
proofs, as do related counterexamples to the full-separability and concurrence
comparisons in the posing source. An elementary bound also establishes its large-qubit limit
without the conjectured shape.
\end{abstract}

\section{The conjecture and its three-qubit instance}
Consiglio, Apollaro and Wie\'sniak call the proposed form of the closest
separable state a conjecture in Appendix A.3 of \cite{CAW}, immediately
before Eq.~(A11). Their closest separable state is closest among
\emph{fully} separable states, as specified in Sections I--II, not among
positive-partial-transpose states. The same conjecture occurs in
\cite[Appendix C.3, Eq.~(C.12)]{Thesis}. We use the source's convention that
``distance'' means the squared Hilbert--Schmidt norm:
\[
 \D(\rho,\sigma)=\Tr((\rho-\sigma)^2),\qquad
 \E(\rho)=\min_{\sigma\in\Sep_n}\D(\rho,\sigma),
\]
where $\Sep_n$ is the convex hull of pure $n$-qubit product projectors.
This hull is compact in finite dimension, and strict convexity of the squared
norm gives a unique closest state. Products of local mixed density matrices
give the same hull, by expanding their spectral decompositions.

Order the three-qubit basis as $000,001,\ldots,111$, with the first qubit
most significant. On the low branch of Eqs.~(A1)--(A2), put
\begin{equation}\label{eq:rho}
 \rho=\diag(1/5,1/5,1/5,1/5,0,0,0,1/5)
       +\frac{11}{112}\bigl(\ket{000}\bra{111}+\ket{111}\bra{000}\bigr).
\end{equation}
It is a density matrix: its corner block has eigenvalues
$1/5\pm11/112$, and the other diagonal entries are nonnegative.
For $n=3$, the family in Eq.~(A11) is
\begin{equation}\label{eq:shape}
 S(a,b,d)=\diag\left(\frac a2,\frac b3,\frac b3,\frac b3,
  \frac{1-a-b}{3},\frac{1-a-b}{3},\frac{1-a-b}{3},\frac a2\right)
  +d\ket{000}\bra{111}+\bar d\ket{111}\bra{000},
\end{equation}
with $a,b\ge0$, $a+b\le1$, and $|d|\le a/2$.
Full separability is an additional requirement.

\begin{theorem}\label{thm:main}
For the state \eqref{eq:rho},
\begin{equation}\label{eq:comparison}
 \E(\rho)\le\frac{417965512477}{90316800000000}
 <\min_{S(a,b,d)\in\Sep_3}\D(\rho,S(a,b,d))
 =\frac{109}{23520}.
\end{equation}
Thus Eqs.~(A8), (A11), and their common form in Eq.~(16) of \cite{CAW}
do not contain a closest fully separable state at this instance.
\end{theorem}
The rational competitor and separability inequality below are formalized
in \cite{Lean}. The proof here also identifies a state attaining the
restricted minimum. Matrix-element separability inequalities provide the
general setting for this argument \cite{GS,CHK}; the particular inequality
needed here is derived directly.

\section{A supporting inequality and the restricted minimum}
Write $p_{xyz}=\sigma_{xyz,xyz}$ and $z=\sigma_{000,111}$.
Every fully separable three-qubit state satisfies
\begin{equation}\label{eq:witness}
 \Rea z\le
 \frac15\left(\frac{p_{001}}6+\frac{3p_{110}}2\right)
 +\frac45\left(\frac{p_{000}}{12}+\frac{p_{011}}{12}
                +\frac{3p_{101}}4+\frac{3p_{110}}4\right).
\end{equation}
To prove this, consider a product of local density matrices with diagonals
$(\alpha_j,\beta_j)$, $j=1,2,3$. Positivity of each $2\times2$ matrix gives
$|z|\le\sqrt{\prod_j\alpha_j\beta_j}$. The weighted arithmetic--geometric
mean inequality gives
\[
 |z|\le\sqrt{p_{001}p_{110}}
       \le p_{001}/6+3p_{110}/2.
\]
For a second bound, first use
$\sqrt{\alpha_1\beta_1}\le\alpha_1/6+3\beta_1/2$.
Multiply by $R=\sqrt{\alpha_2\beta_2\alpha_3\beta_3}$, and use
\[
 R\le(\alpha_2\alpha_3+\beta_2\beta_3)/2,
 \qquad R\le(\alpha_2\beta_3+\beta_2\alpha_3)/2
\]
in the terms with coefficients $\alpha_1/6$ and $3\beta_1/2$, respectively.
This gives the second parenthesis in \eqref{eq:witness}.
Both bounds are linear in matrix entries after taking $\Rea z$.
They therefore pass to convex mixtures; taking weights $1/5$ and $4/5$
proves \eqref{eq:witness}.

For \eqref{eq:shape}, set
\[
 A=a/2,\quad B=b/3,\quad C=(1-a-b)/3,\quad x=\Rea d,\quad y=\Ima d.
\]
Then \eqref{eq:witness} says
\begin{equation}\label{eq:halfspace}
 h:=A/15+B/10+3C/2-x=\frac12-\frac7{15}(a+b)-x\ge0.
\end{equation}
Expanding the distance and completing squares gives the exact identity
\begin{align}\label{eq:squares}
 \D(\rho,S)
 &=2(1/5-A)^2+3(1/5-B)^2+3C^2+2(11/112-x)^2+2y^2\notag\\
 &=\frac{109}{23520}+2(A-3/16)^2+3(B-3/16)^2
       +3(C-1/48)^2\notag\\
 &\hspace{18mm}+2(x-1/16)^2+2y^2+\frac h7.
\end{align}
This proves the lower bound for every fully separable matrix of the
proposed shape, including complex $d$ and all boundary parameters.

Here is an explicit separable state attaining equality.
For $0\le t\le1$ and $r\in\mathbb Z/3\mathbb Z$, define
\[
 v(t,r)=\sqrt t\ket0+\omega^r\sqrt{1-t}\ket1,
 \qquad\omega=e^{2\pi i/3},
\]
and the product-state average
\begin{equation}\label{eq:twirl}
 T(t_1,t_2,t_3)=\frac19\sum_{r,s=0}^2
  \proj{v(t_1,r)\otimes v(t_2,s)\otimes v(t_3,-r-s)}.
\end{equation}
The vectors are normalized, so $T$ is fully separable. Its diagonal entry
at a word $q$ is $\prod_j t_j^{1-q_j}(1-t_j)^{q_j}$, and its corner
coherence is $\prod_j\sqrt{t_j(1-t_j)}$.
Every other off-diagonal entry is zero: averaging the two characters forces
$q_1-q'_1=q_3-q'_3$ and $q_2-q'_2=q_3-q'_3$ modulo $3$.
Each difference lies in $\{-1,0,1\}$, so they must be equal as integers.
This leaves the diagonal and the two opposite corners only.
Consequently
\begin{equation}\label{eq:sstar}
 S_*:=\frac56T(9/10,1/2,1/2)+\frac16\proj{111}
     =S(3/8,9/16,1/16)
\end{equation}
is fully separable and makes every nonconstant term in
\eqref{eq:squares} zero. This proves the equality in \eqref{eq:comparison}.

\section{A closer product-state mixture}
Now change the last two parameters in \eqref{eq:sstar}:
\begin{equation}\label{eq:tau}
 \tau=\frac56T(9/10,199/400,199/400)+\frac16\proj{111}.
\end{equation}
This is a convex mixture of ten pure product projectors, nine with weight
$5/54$ and one with weight $1/6$. In binary order, its diagonal is
\[
 \frac1{1920000}(356409,359991,359991,363609,
                 39601,39999,39999,360401),
\]
and its only off-diagonal entries are
$\tau_{000,111}=\tau_{111,000}=119997/1920000$.
For example, the corner coherence of the unweighted average is
\[
 \sqrt{\tfrac9{10}\tfrac1{10}}\,
 \left(\sqrt{\tfrac{199}{400}\tfrac{201}{400}}\right)^2
 =\frac{119997}{1600000}.
\]
Entrywise squaring gives
\begin{equation}\label{eq:gap}
 \D(\rho,\tau)=\frac{417965512477}{90316800000000},\qquad
 \frac{109}{23520}-\D(\rho,\tau)
   =\frac{84926789}{12902400000000}>0.
\end{equation}
Together with \eqref{eq:squares}, this proves Theorem~\ref{thm:main}.
The state $\tau$ still respects interchange of qubits 2 and 3; it violates
the additional population equalities in \eqref{eq:shape}.
In particular, it remains diagonal apart from the corner coherence.
The obstruction is to the population restrictions, not to this sparse support.

Target invariance alone does not justify those restrictions. The unitary
swapping only $\ket{000}$ and $\ket{111}$ fixes $\rho$ but sends the product
$(\ket{000}+\ket{001})/\sqrt2$ to
$(\ket{111}+\ket{001})/\sqrt2$, which has Schmidt rank two across the
first-qubit cut. It does not preserve the feasible set.
By contrast, permutations of qubits $2,\ldots,n$ and correlated local
diagonal phases do preserve full separability and fix a low-branch target.
Averaging over these symmetries cannot increase its squared distance.
In particular, the finite local-phase average used in \eqref{eq:twirl}
extends to $n$ qubits, with the last phase the negative sum of the first
$n-1$ phases. The same character argument retains exactly the diagonal and
corner coherences. Thus a closest state can always be chosen with this
sparse support and with populations invariant under those qubit permutations;
the unjustified additional equalities are what must be removed.

\section{Related statements in the posing source}
\subsection{The zero-negativity condition is not full separability}
\begin{corollary}\label{cor:ppt}
The condition in \cite[Eq.~(A10)]{CAW} is insufficient for full separability
of the family (A8). The state
\[
 \eta=S(5/16,5/8,33/512)
\]
is positive under partial transpose across every one-versus-two cut and is
separable across each such cut, but is not fully separable.
\end{corollary}
\begin{proof}
Its diagonal entries are $5/32$ at the endpoints, $5/24$ at the three
upper interior words, and $1/48$ at the lower interior words.
The original corner block is positive since $33/512<5/32$.
Each partial transpose has one potentially nonpositive $2\times2$ block,
whose determinant is
\[
 \frac5{24}\frac1{48}-\left(\frac{33}{512}\right)^2
 =\frac{439}{2359296}>0.
\]
For any fixed cut, the coherent block acts on a $2\otimes2$ subspace
spanned on the other side by $\ket{00}$ and $\ket{11}$.
It is separable by the $2\otimes2$ positive-partial-transpose criterion
\cite{HHH}; the remaining diagonal terms are product projectors across that
cut. Nevertheless, \eqref{eq:halfspace} requires
$\Rea d\le1/16$, whereas $33/512=1/16+1/512$.
\end{proof}
This example concerns the sufficiency of the parameter condition, not
whether a particular solution of the source's restricted optimization is
bound entangled. It also excludes full-separability sufficiency for (A12)
in general, since (A10) is its three-qubit specialization.

\subsection{Concurrence does not determine the squared distance}
The two-qubit comparisons in \cite[Eqs.~(17)--(18)]{CAW} are false.
Let $R$ have diagonal $(1/3,1/3,0,1/3)$ and real corner coherence $1/4$.
Its concurrence is $C(R)=1/2$: the square roots in the concurrence formula
\cite{Wootters} are $7/12,1/12,0,0$.
A fully separable competitor $Q$ has diagonal
\[
 (7/24,1/3,1/12,7/24)
\]
and corner coherence $1/6$.
Indeed, with
\[
 T_2(t_1,t_2)=\frac13\sum_{r=0}^2
 \proj{v(t_1,r)\otimes v(t_2,-r)},
\]
one has $Q=\frac34T_2(2/3,1/3)+\frac18\proj{00}+\frac18\proj{11}$.
Thus $\E(R)\le\D(R,Q)=7/288$.

For $0\le c\le1$, put $P=\proj{\Phi^+}$, where
$\ket{\Phi^+}=(\ket{00}+\ket{11})/\sqrt2$, and set
\[
 J_c=\frac{1+c}{2}P+\frac{1-c}{6}(I-P).
\]
Its Bell-basis eigenvalues give $C(J_c)=c$, and
\begin{equation}\label{eq:isotropic}
 \E(J_c)=c^2/3.
\end{equation}
For completeness, any separable $\sigma$ obeys $\Tr(P\sigma)\le1/2$.
The traceless matrix $W=P-I/4$ has squared norm $3/4$, so Cauchy--Schwarz
gives $\D(J_c,\sigma)\ge c^2/3$.
Equality is attained by $J_0=I/6+P/3$, the equal mixture of the six product
projectors $\proj{u\otimes\bar u}$, with
\[
 u\in\{\ket0,\ket1,(\ket0\pm\ket1)/\sqrt2,
                       (\ket0\pm i\ket1)/\sqrt2\}.
\]
At $c=1/2$, the concurrences of $R$ and $J_c$ agree but
$\E(R)\le7/288<1/12=\E(J_{1/2})$, disproving (17).
At $c=2/5$, $C(R)>C(J_c)$ but
$\E(R)\le7/288<4/75=\E(J_{2/5})$, disproving (18).

The multipartite failures reported in \cite[Eqs.~(19)--(20)]{CAW} also
have exact examples within three-qubit $X$-states. Corollary~\ref{cor:ppt}
makes $C_{\mathrm{GME}}(\eta)=0$, while $\E(\eta)>0$.
Comparison with a diagonal state disproves equivalence of equality.
For a strict reversal, write \eqref{eq:witness} as
$\Tr(H\sigma)\le0$, with $H_{000,111}=H_{111,000}=1/2$ and diagonal
\[
 (-1/15,-1/30,0,-1/15,0,-3/5,-9/10,0).
\]
Then $\Tr(H\eta)=1/512$ and $\Tr(H^2)=42/25$, whence
\[
 \E(\eta)\ge\frac{25}{11010048}.
\]
The low-branch three-qubit state $\rho_{\epsilon}$ with corner coherence
$\epsilon=1/2048$ has
$C_{\mathrm{GME}}(\rho_{\epsilon})=2\epsilon>0$
by the $X$-state concurrence formula \cite{Rafsanjani}, but its diagonal
part gives
\[
 \E(\rho_{\epsilon})\le2\epsilon^2=\frac1{2097152}
 <\frac{25}{11010048}\le\E(\eta).
\]
Thus the order reversal in (20) holds even with one state separable across
every bipartition.

\subsection{The upper bound and large-qubit limit}
\begin{corollary}\label{cor:limit}
Let $X_n$ be any $n$-qubit density matrix whose only off-diagonal entries
are the opposite corners $\gamma_n,\bar\gamma_n$. Then
\begin{equation}\label{eq:limit}
 2\bigl(\max\{|\gamma_n|-2^{-n},0\}\bigr)^2
 \le\E(X_n)\le2|\gamma_n|^2.
\end{equation}
Consequently, if $|\gamma_n|\to g$, then $\E(X_n)\to2g^2$.
\end{corollary}
\begin{proof}
The diagonal of $X_n$ is fully separable and gives the upper bound.
Every local qubit density matrix has off-diagonal magnitude at most $1/2$,
because $|q_{01}|^2\le q_{00}q_{11}\le1/4$.
Hence each product, and then every fully separable mixture, has corner
coherence at most $2^{-n}$. The two corner entries alone contribute at least
the left side of \eqref{eq:limit} to the squared norm.
The last assertion follows by squeezing.
\end{proof}
This proves the upper bound and the fixed-coherence limit in
\cite[Eq.~(A14)]{CAW} for normalized X-MEMS, independently of (A11).
It also covers convergent coherences and arbitrary diagonal populations.

\section{Coverage and verification scope}
\begin{center}
\small
\begin{tabular}{@{}p{0.23\linewidth}p{0.69\linewidth}@{}}
\toprule
Source item in \cite{CAW} & Coverage\\
\midrule
(16), (A8), (A11) & Refuted at $n=3$, $\gamma=11/112$ by Theorem~\ref{thm:main}; the restricted minimum is exact.\\
(A10), (A12) & Necessary PPT bounds; not sufficient for full separability, by Corollary~\ref{cor:ppt}.\\
(17), (18) & Refuted by $R$ and $J_c$, using $c=1/2$ and $2/5$.\\
(19), (20) & Reported failures proved with $\eta$, a diagonal state, and $\rho_{1/2048}$.\\
(A14), Fig.~6 bound & Proved by Corollary~\ref{cor:limit}, independently of the shape conjecture.\\
(A13), GHZ endpoint & Prior formula in \cite[Eq.~(9)]{Pandya}; not a consequence of the refuted shape reduction.\\
(A5)--(A7), Fig.~4 & Two-qubit formulas not reassessed here; the concurrence counterexamples do not invalidate those formulas.\\
Figs.~5--6, quartic rule & The shape reduction is invalid as a general certificate; pointwise numerical values and a corrected general formula are not determined here.\\
Section V, $k$-CSS and algorithm & The general closest-$k$-separable-state algorithm and its performance are not covered.\\
\bottomrule
\end{tabular}
\end{center}
The main counterexample corresponds to the Lean declaration
\texttt{result : $\neg$ claim} in the namespace
\begin{quote}\small\ttfamily
D5.S3.Quantum.Entanglement.GeometricMeasure.\\
XMemsClosestSeparableShapeRefutation
\end{quote}
The \href{https://github.com/the-omega-institute/trureturing/blob/83b70a1d0d6a58bb32521c741805d845d401c5f8/D5/S3/Quantum/Entanglement/GeometricMeasure/XMemsClosestSeparableShapeRefutation.lean}{pinned Lean file}
defines the low-branch target, the full parameter family and the
finite-convex-mixture notion of full separability. Its declaration proves the
negation of the universal closest-state shape assertion by the instance above.
The restricted-minimum attainment, the symmetry reduction and the consequences
in Section~4 have written proofs in this note. The accompanying
\texttt{verify.py} checks the rational matrices, character support at all
64 three-qubit matrix positions, the four-variable square-completion identity,
and the finite constants in the corollaries.

The actual global minimum at \eqref{eq:rho} and an explicit parameterization
of closest states for general finite $n$ remain to be determined.

\paragraph{License.}
This note and its accompanying original source and verification code are
licensed under Creative Commons Attribution 4.0 International (CC BY 4.0),
\url{https://creativecommons.org/licenses/by/4.0/}.

\begin{thebibliography}{10}
\bibitem{CAW}
M.~Consiglio, T.~J.~G.~Apollaro and M.~Wie\'sniak,
\emph{Variational approach to the quantum separability problem},
Phys. Rev. A \textbf{106}, 062413 (2022),
\href{https://doi.org/10.1103/PhysRevA.106.062413}{doi:10.1103/PhysRevA.106.062413}.
Posing text: \href{https://arxiv.org/abs/2209.01430v2}{arXiv:2209.01430v2},
Appendix A.3, Eq.~(A11); Sections I--II and IV.

\bibitem{Thesis}
M.~Consiglio, \emph{Variational Quantum Algorithms for Many-Body Systems},
Ph.D. thesis, University of Malta (2025),
\href{https://arxiv.org/abs/2502.11985v1}{arXiv:2502.11985v1},
Appendix C.3, Eq.~(C.12), printed pp.~118--119.

\bibitem{Lean}
The Omega Institute, \emph{XMemsClosestSeparableShapeRefutation.lean},
\texttt{trureturing}, revision
\texttt{83b70a1d0d6a58bb32521c741805d845d401c5f8}.
\href{https://github.com/the-omega-institute/trureturing/blob/83b70a1d0d6a58bb32521c741805d845d401c5f8/D5/S3/Quantum/Entanglement/GeometricMeasure/XMemsClosestSeparableShapeRefutation.lean}{Source file and declaration \texttt{result}}.

\bibitem{GS}
O.~G\"uhne and M.~Seevinck,
\emph{Separability criteria for genuine multiparticle entanglement},
New J. Phys. \textbf{12}, 053002 (2010),
\href{https://doi.org/10.1088/1367-2630/12/5/053002}{doi:10.1088/1367-2630/12/5/053002};
\href{https://arxiv.org/abs/0905.1349}{arXiv:0905.1349}, Observation 4.

\bibitem{CHK}
L.~Chen, K.~H.~Han and S.-H.~Kye,
\emph{Separability criterion for three-qubit states with a four dimensional norm},
J. Phys. A: Math. Theor. \textbf{50}, 345303 (2017),
\href{https://doi.org/10.1088/1751-8121/aa7f9a}{doi:10.1088/1751-8121/aa7f9a};
\href{https://arxiv.org/abs/1705.05557v2}{arXiv:1705.05557v2}, Theorem 2.1.

\bibitem{HHH}
M.~Horodecki, P.~Horodecki and R.~Horodecki,
\emph{Separability of mixed states: necessary and sufficient conditions},
Phys. Lett. A \textbf{223}, 1--8 (1996),
\href{https://doi.org/10.1016/S0375-9601(96)00706-2}{doi:10.1016/S0375-9601(96)00706-2};
\href{https://arxiv.org/abs/quant-ph/9605038}{arXiv:quant-ph/9605038}.

\bibitem{Wootters}
W.~K.~Wootters, \emph{Entanglement of formation of an arbitrary state of two qubits},
Phys. Rev. Lett. \textbf{80}, 2245--2248 (1998),
\href{https://doi.org/10.1103/PhysRevLett.80.2245}{doi:10.1103/PhysRevLett.80.2245};
\href{https://arxiv.org/abs/quant-ph/9709029}{arXiv:quant-ph/9709029}, Eqs.~(9)--(10).

\bibitem{Rafsanjani}
S.~M.~Hashemi Rafsanjani, M.~Huber, C.~J.~Broadbent and J.~H.~Eberly,
\emph{Genuinely multipartite concurrence of $N$-qubit $X$ matrices},
Phys. Rev. A \textbf{86}, 062303 (2012),
\href{https://doi.org/10.1103/PhysRevA.86.062303}{doi:10.1103/PhysRevA.86.062303};
\href{https://arxiv.org/abs/1208.2706v2}{arXiv:1208.2706v2}, Eq.~(5).

\bibitem{Pandya}
P.~Pandya, O.~Sakarya and M.~Wie\'sniak,
\emph{Hilbert--Schmidt distance and entanglement witnessing},
Phys. Rev. A \textbf{102}, 012409 (2020),
\href{https://doi.org/10.1103/PhysRevA.102.012409}{doi:10.1103/PhysRevA.102.012409};
\href{https://arxiv.org/abs/1811.06599}{arXiv:1811.06599}, Eq.~(9).
\end{thebibliography}
\end{document}
